We have presented a systematic study of machine-learning detection of
controlled beyond-GR deviations in gravitational-wave signals, using
synthetic and real LIGO noise. A hybrid classifier combining a 1D CNN
with hand-crafted waveform statistics detects amplitude, phase, and
frequency modulations with high accuracy at moderate to large
deviation strengths. Using the real GW150914 strain as a template we
measure a detectability curve with threshold $\beta \approx 0.25$,
where the value of $\beta$ is specific to the quadratic-in-time
modulation form adopted here. Detection probability rises smoothly
from chance at $\beta \leq 0.2$ to perfect classification at
$\beta \geq 0.5$. The classifier generalizes to unseen deviation
types and to a wide range of binary mass ratios. At small $\beta$
values typical of current parametrized-GR constraints, our method
fails; this negative result quantifies the limits of ML-only
beyond-GR searches in real detector noise. These results establish
both the promise and the current limits of machine learning for
beyond-GR tests with gravitational waves. Closing the sensitivity gap
to matched filtering---potentially by combining the two approaches---
is a natural direction for future research.
